\begin{mistake}{2026-05-25}{02}

\begin{problemcontent}
计算数列极限 \(I = \lim_{n \to \infty} n \tan(\pi \sqrt{n^2 + 1})\)
\end{problemcontent}

\begin{solutioncontent}
当 \(n \to \infty\) 时，\(\sqrt{n^2+1} \approx n\)，由此构造整数周期 \(n\pi\)，利用正切函数的周期性 \(\tan(x + n\pi) = \tan x\) 将自变量由无穷大转化为无穷小：
\[
\begin{aligned}
I &= \lim_{n \to \infty} n \tan(\pi \sqrt{n^2 + 1}) \\
&= \lim_{n \to \infty} n \tan(\pi \sqrt{n^2 + 1} - n\pi) \\
&= \lim_{n \to \infty} n \tan\Big( \pi(\sqrt{n^2 + 1} - n) \Big) 
\end{aligned}
\]
对括号内进行分子有理化（根式差变形）：
\[
\begin{aligned}
I &= \lim_{n \to \infty} n \tan\left( \pi \cdot \frac{(n^2 + 1) - n^2}{\sqrt{n^2 + 1} + n} \right) \\
&= \lim_{n \to \infty} n \tan\left( \frac{\pi}{\sqrt{n^2 + 1} + n} \right)
\end{aligned}
\]
当 \(n \to \infty\) 时，\(\frac{\pi}{\sqrt{n^2 + 1} + n} \to 0\)，使用等价无穷小 \(\tan x \sim x\)：
\[
\begin{aligned}
I &= \lim_{n \to \infty} n \cdot \frac{\pi}{\sqrt{n^2 + 1} + n} \\
&= \lim_{n \to \infty} \frac{\pi}{\sqrt{1 + \frac{1}{n^2}} + 1} \\
&= \frac{\pi}{2}
\end{aligned}
\]
\end{solutioncontent}

\mistakeknowledge{
\begin{itemize}
  \item \textbf{核心公式/定理}：\(\tan(x \pm n\pi) = \tan x\)。对于三角函数内趋于 \(\infty\) 的自变量，常通过剥离周期构造趋于 \(0\) 的无穷小。
  \item \textbf{易错点}：切忌直接将 \(\sqrt{n^2+1}\) 近似为 \(n\) 从而得出 \(0\) 的错解，函数内部不可直接粗暴舍弃低阶项。
  \item \textbf{关联知识}：若原题换成 \(\sin(\pi \sqrt{n^2 + 1})\)，则减去 \(n\pi\) 会变为 \(\sin(\pi \sqrt{n^2+1} - n\pi + n\pi) = (-1)^n \sin(\pi(\sqrt{n^2+1}-n))\)，需注意交错符号。
\end{itemize}}

\end{mistake}
